% ZQE objective in likelihood form
% Notation follows main.tex: \bt = decoder params, q_{\Z\mid\Y;\bt} = surrogate,
% f_{\Y\mid\Z;\bt} = decoder density, f_{\Y;\bt} = marginal.

\subsection*{$Z_q$ objective in log-likelihood form}

The $Z_q$ estimating equations $\Psi_n(\bt) = 0$ are the first-order conditions of the
scalar objective
%
\begin{equation}
  \mathcal{L}_{\mathrm{ZQE}}(\bt)
  \;:=\;
  \underbrace{
    \frac{1}{n}\sum_{i=1}^n
    \mathbb{E}_{q_{\Z\mid\Y;\bt}(\cdot\mid\y_i)}
    \!\bigl[\log f_{\Y\mid\Z;\bt}(\y_i\mid\Z)\bigr]
  }_{\text{surrogate expected log-lik on real data}}
  \;-\;
  \underbrace{
    \mathbb{E}_{f_{\Y;\bt}(\y)}
    \mathbb{E}_{q_{\Z\mid\Y;\bt}(\cdot\mid\y)}
    \!\bigl[\log f_{\Y\mid\Z;\bt}(\y\mid\Z)\bigr]
  }_{\text{same quantity, centred under model}}.
  \label{eq:L_ZQE}
\end{equation}
%
The first term is the surrogate-posterior expected conditional log-likelihood evaluated
on the observed data; the second term is its expectation under the generative model
$f_{\Y;\bt}$, and plays the role of a \emph{self-centering} baseline.

\paragraph{Why this is a valid scalar objective.}
For exponential-family decoders, differentiating $\mathcal{L}_{\mathrm{ZQE}}$ with
respect to $\bt$ (treating the dependence of $q_{\Z\mid\Y;\bt}$ on $\bt$ as fixed,
i.e.\ holding the surrogate at its current value) recovers exactly the $Z_q$ estimating
map $\Psi_n(\bt)$ from equation~\eqref{eq:Psi_n_general}.
More precisely, holding the surrogate fixed and differentiating under the integral sign,
%
\begin{align*}
  \nabla_\bt \mathcal{L}_{\mathrm{ZQE}}(\bt)
  &=
  \frac{1}{n}\sum_{i=1}^n
  \mathbb{E}_{q_{\Z\mid\Y;\bt}(\cdot\mid\y_i)}
  \!\bigl[\nabla_\bt \log f_{\Y\mid\Z;\bt}(\y_i\mid\Z)\bigr]
  -
  \mathbb{E}_{f_{\Y;\bt}}
  \mathbb{E}_{q_{\Z\mid\Y;\bt}(\cdot\mid\cdot)}
  \!\bigl[\nabla_\bt \log f_{\Y\mid\Z;\bt}(\Y\mid\Z)\bigr]
  \\[4pt]
  &=\;
  \Psi_n(\bt).
\end{align*}
%
Hence minimising $\mathcal{L}_{\mathrm{ZQE}}$ over $\bt$ is equivalent to solving the
$Z_q$ estimating equations in the score direction.

\paragraph{Practical computation (fixed-seed Monte Carlo).}
Both expectations are approximated by Monte Carlo.
For the first term, draw $z_m \sim q_{\Z\mid\Y;\bt}(\cdot\mid\y_i)$ for $m=1,\dots,M$
via reparameterisation:
%
\[
  z_m = \mu_q(\y_i;\bt) + \sigma_q(\y_i;\bt)\,\odot\,\varepsilon_m,
  \qquad \varepsilon_m \sim \mathcal{N}(0,I).
\]
%
For the second term, draw $(\tilde\y_m, \tilde z_m)$ from the joint generative model:
$\tilde z_m \sim f_\Z$, $\tilde\y_m \sim f_{\Y\mid\Z;\bt}(\cdot\mid\tilde z_m)$.
Fixing the seeds $\{\varepsilon_m\}$ and $\{(\tilde z_m, \tilde\y_m)\}$ before
optimising makes the objective deterministic in $\bt$, so that a quasi-Newton method
(L-BFGS) is applicable with valid curvature estimates throughout the solve.
The reparameterisation ensures that $\nabla_\bt$ still flows through $z_m(\bt)$ even
though the noise $\varepsilon_m$ is fixed.

\paragraph{Connection to the ELBO.}
Dropping the centering term recovers the ELBO:
%
\[
  \mathcal{L}_{\mathrm{ELBO}}(\bt)
  =
  \frac{1}{n}\sum_{i=1}^n
  \mathbb{E}_{q_{\Z\mid\Y;\bt}(\cdot\mid\y_i)}
  \!\bigl[\log f_{\Y\mid\Z;\bt}(\y_i\mid\Z)\bigr]
  -
  \mathrm{KL}\!\bigl(q_{\Z\mid\Y;\bt}(\cdot\mid\y_i)\;\|\;f_\Z\bigr).
\]
%
The $Z_q$ objective differs in two respects: (i) it replaces the KL penalty by the
model-expectation centering term, and (ii) it does not penalise the surrogate for
deviating from the prior. As a result, the population maximum of
$\mathcal{L}_{\mathrm{ZQE}}$ at any surrogate always coincides with $\bt_0$, whereas
the population ELBO maximum generally does not when $q \neq f_{\Z\mid\Y;\bt}$.
